\section{Threats to validity and limitations}
\label{sec:threats}

\subsection{Threats to the validity of the measurements}

\paragraph{The baselines.} The playbooks compare against Mesa~2.3.2, the version the ported
research model was written for, and the current Mesa~3.5.1; Mesa~2.3 runs here on CPython
3.15, which it does not officially support. The Mesa~3.5 and mesa-frames implementations were
written for this paper, following each library's documentation and examples, and their
developers might write them faster. mesa-frames was measured from its development branch
(commit \code{ef92326}, polars~2.0, on one thread), on four of the eight workloads; writing
the others needs features it lacks (it has no continuous or network space), and working
around three defects we met (neighbourhoods on a torus, the removal of agents, and
neighbourhoods of radius above one) made its code longer than its examples. AgentPy,
Agents.jl and NetLogo have not been measured, and the speed-ups should not be read as applying
to them.

\paragraph{Framework or algorithm.} The playbook speed-ups measure what a modeller gains by
writing each model in the form the framework offers, not the framework's overhead alone. Part
of a gain can be algorithmic: the Mesa flocking implementation scans every agent for
neighbours, as Mesa's continuous space does, while the \modelly{} one queries a KD-tree; a
Mesa model could call \code{scipy.spatial.cKDTree} itself. The cell-list variants keep Mesa's
data structure but not all of its loops: in Schelling's model most of their gain comes from
counting every resident's similar neighbours at once, which takes three quarters of Mesa's run
time, while Mesa's move to a random empty cell already runs in constant expected time. In
Sugarscape the Mesa implementations and the cell-list variant run the same search algorithm, so
their difference is that of the data structure; the column and mesa-frames variants vectorise
the search as well.

\paragraph{What the timings include.} The timed region is the whole simulation, including the
data collection and the summary statistics of each workload. In the Boltzmann wealth model the
data collector (a Gini coefficient of 10,000 agents every step) takes about 70\% of
the run on Mesa~2.3, two thirds on \modelly{} and about half on Mesa~3.5, so that workload
measures data collection as much as the exchange; the Mesa implementations compute the Gini
coefficient in Python, as Mesa's own example does, and \modelly{}'s with \code{numpy}. The summaries computed after the last step are a
small part of each run (about 2\% of Mesa's Schelling run).

\paragraph{Who wrote the implementations.} Every implementation of every playbook, the Mesa ones
included, was written by the author of \modelly{} with the help of AI coding assistants. Each
Mesa version follows Mesa's documented API and examples, and all of them are in the repository
for inspection; corrections from Mesa developers are welcome.

\paragraph{Measurement.} All timings come from one laptop with a hybrid processor, in one
twelve-minute session, pinned to one performance core (CPU 2, one of the processor's four
4.7~GHz cores; the other twelve are efficiency cores), under the \code{intel\_pstate} driver
with its performance energy preference and turbo enabled, while other work shared the machine
(one-minute load average 1.3 to 5.7 at the start of each workload, 17 at the start of the first,
right after the agreement study). Frequency could not be fixed without administrator rights.
The speed-ups are ratios of medians over ten rounds, the order of the implementations
rotating, which removes most of the drift but not all: the per-round ratios lie within 25\% of
the median for 14 of the 25 non-reference variants and within 38\% for ten more, the widest
being Mesa~3.5 on the insect population ($-44\%$ to $+35\%$; \cref{fig:speedups}). Earlier
sessions, on an efficiency core and under more load, gave ratios up to 30\% different for some
workloads, so the conclusions rest on the order of magnitude of each speed-up, not on its last
digit.

\paragraph{Agreement.} The agreement study (\cref{sec:agreement}) compares chosen summary
statistics, not whole trajectories; six small differences are individually detectable at the
small size, though none survives the correction; at the timed size it uses 50 seeds (200 for
the grid epidemic), whose intervals are wider; and an equivalence margin of half a standard
deviation is a choice: it is a medium effect by the usual convention, the margin is estimated
from the same data, which makes the test slightly liberal, and stricter margins give fewer
equivalences (27 of 67 at $0.3\,s_p$, 19 at $0.25\,s_p$).

\paragraph{Reproducibility.} Continuous integration runs the determinism suite with the
pure-Python core on all three operating systems; the compiled core's equality with it is
established on the development machine only (\cref{tab:determinism}), and the golden digests
cover ten steps of fourteen examples. Running the suite on every published wheel, for more
steps, is planned.

\paragraph{A young project.} \modelly{} was first released on 24 October 2026, by one developer;
its seven releases to date span one day, the policy of \cref{sec:regressions} has been applied
to two of them (1.3.1 and 1.3.2), and release 1.3.0 shipped two regressions that 1.3.1 fixed. The research workflow was added in 1.3
and its public names are marked provisional before 2.0.

\subsection{Limitations of the software}
\label{sec:limitations}

\begin{itemize}[leftmargin=*,itemsep=2pt]
  \item \textbf{Speed comes from the model's form.} A model written agent by agent gains only
  where the framework's own calls take its time (none to 4.5 times, \cref{sec:portability}); the
  large gains of \cref{sec:playbooks} require phases that can be
  expressed over the whole population. Sequential decisions where each agent depends on the
  previous one's move remain per-agent loops.
  \item \textbf{One machine per run.} Experiments distribute independent runs across processes
  and clusters, but a single run executes in one process. Spatially partitioned runs across
  machines, a GPU backend for column phases, and differentiable models are planned and not yet
  available.
  \item \textbf{Checkpoints and closures.} Checkpoints pickle the model; models that hold
  lambdas, nested functions or classes created at run time cannot be saved, and the framework
  refuses them with an error naming the object. Checkpoints load only in the \modelly{} minor
  version that wrote them.
\end{itemize}
